% Generated by logic/gen/gen_paper_latex.py -- do not edit by hand.
%
% Jinja runs here with non-default delimiters (block, variable and comment
% markers all use angle-bracket forms) so that LaTeX's own braces pass through
% untouched. See report_utils.render_template for the exact configuration.
%
% Context: label, caption, column_spec, first_header, headers[], rows[] of
% {label, cells[]}, and an optional note rendered as a table footnote.
\begin{table}[htbp]
\centering
\caption{Mandatory-selection strategies at the 30-day horizon, ordered by observed collection efficiency. Efficiency falls monotonically as the rule grows more conservative; overflow risk does not fall with it, being essentially flat across the three middle strategies.}\label{tab:strategies30}
\begin{tabular}{lrrrrrr}
\toprule
Selection strategy & kg/km & Overflows & kg lost & km & Collections & Time (s) \\
\midrule
Last-Minute (CF90) & \textbf{7.74} & 28.6 & 163.4 & \textbf{3,725} & 2,141 & 1,144 \\
Look-Ahead & 6.60 & 6.3 & 30.0 & 4,327 & 2,502 & \textbf{1,121} \\
Last-Minute (CF70) & 6.23 & 7.7 & 38.2 & 4,646 & 2,721 & 1,446 \\
Service-Level (SL1) & 5.65 & 6.3 & 25.0 & 5,088 & 3,037 & 1,712 \\
Service-Level (SL2) & 4.21 & \textbf{1.4} & \textbf{2.0} & 6,696 & 4,036 & 2,263 \\
\bottomrule
\end{tabular}

\vspace{0.4em}
\begin{minipage}{\linewidth}
\footnotesize CF70/CF90 are Last-Minute critical-fill thresholds of 70\% and 90\%; SL1/SL2 are Service-Level projection horizons of one and two days.
\end{minipage}

\end{table}
